\section{Add inflation while retaining the collector}
Use \path{packages/033-png/construction/png-inflater-construction/attach.py} with your current completed
collector program as \texttt{--source} and a fresh \texttt{--output}. It replaces the
previous supplied temporary wrapper and appends unfinished inflater regions. It keeps
your loader and collector instructions. If you changed that wrapper, reconcile it by
hand; the helper refuses to discard an unrecognized implementation.

The new wrapper retains the collector's private descriptor until it has passed the
joined stream address and count to \texttt{png\_inflate}. The row-workspace address
comes from decoder context offset 32. The predicted exact output count comes from
collector descriptor offset 32; those equal offsets belong to different records and
have different meanings. The wrapper checks that this count fits the actual row capacity
at context offset 40. Stop at \texttt{inflate\_call} to inspect the four arguments and
at \texttt{inflate\_return} to inspect the result. Final context outputs and framebuffer
words remain unpublished at this stage.

The next sixteen edits develop a complete stored-block path. Fixed and dynamic Huffman
paths retain explicit development exits until their later construction steps. This
intermediate program is therefore not a completed inflater submission. It can return
success for a valid stored stream and development status for a Huffman block that it
has not yet attempted to decode. The reserved block type is already a failure.

The supplied 2048-byte frame leaves room for later call-private code tables. It saves
\texttt{\$ra}, the eight saved registers and \texttt{\$fp} at offsets 64 through 136.
The frame pointer identifies this call's private workspace. A scratch doubleword at
offset zero retains a length or checksum across helper calls; offset 56 retains the
advertised history-window size. The illustrated state uses \texttt{\$s0} for the next
compressed byte, \texttt{\$s1} for the trailer's first byte, \texttt{\$s2}/\texttt{\$s3}
for reservoir bits and their count, \texttt{\$s4}/\texttt{\$s5} for output cursor and
expected end, \texttt{\$s6} for the current final flag, and \texttt{\$s7} for output
start. These are checkpoint roles, not requirements on an equivalent final solution.

Internal unfinished helpers unwind the owning inflater frame and return development
status four with count zero. Call them through \texttt{png\_inflate} while constructing
this sequence. A helper is not safe to enter on its own before its body is complete.
The preparation helper retains instruction edits and selects the named input. Its
fixture helper can write a PNG for the supplied program or a raw zlib stream for a
direct-entry observation. Keep those two file formats distinct when choosing arguments.
